\documentclass[11pt]{article}

\setlength{\parindent}{0pt}
\usepackage{xltxtra}
\usepackage{hyperref}
\hypersetup{hidelinks}
\usepackage{url}
\urlstyle{tt}
\usepackage{xcolor}
\definecolor{CVBlue}{RGB}{23,110,191}
\usepackage{calc}
\usepackage{graphicx}
\usepackage{tikz}
\usetikzlibrary{calc}
\usepackage{fontspec}
\usepackage{xeCJK}
\usepackage{enumitem}
\CJKsetecglue{} 

%% 主文档字体设置
\setmainfont[
    Path = fonts/Main/,
    Extension = .otf,
    BoldFont = texgyretermes-bold.otf,
]{texgyretermes-regular.otf}

% 中文字体设置
\setCJKmainfont[
    Path = fonts/hansans/,
    Extension = .ttf,
    BoldFont = NotoSansSC-Bold.ttf,
]{NotoSansSC-Regular.otf}

\usepackage{relsize}
\usepackage{xspace}
\usepackage{fontawesome} 

% A4纸，上下左右边距
\usepackage[
    a4paper,
    left=1.2cm,
    right=1.2cm,
    top=1.5cm,
    bottom=1cm,
    nohead
]{geometry}

\renewcommand{\baselinestretch}{1.3} % 稍微收紧行间距以容纳更多干货

\usepackage{titlesec}
\usepackage{enumitem}
\setlist{noitemsep} 
\setlist[itemize]{topsep=0.15em, leftmargin=*}
\setlist[enumerate]{topsep=0.15em, leftmargin=*}

\newlength{\interProjectSpacing}
\setlength{\interProjectSpacing}{0.7em} 
\newcommand{\projectsep}{\vspace{\interProjectSpacing}}

\newlength{\intraProjectTitleSep}
\setlength{\intraProjectTitleSep}{0.3em} 
\newcommand{\titlebreak}{\\[\intraProjectTitleSep]}

\titleformat{\section}         
{\large\bfseries\raggedright} 
{}{0em}                      
{}                           
[{\color{CVBlue}\titlerule}]  
\titlespacing*{\section}{0cm}{*1.2}{*0.8}

\begin{document}
\pagenumbering{gobble}

\centerline{\LARGE\bfseries{林十辉}} 

\centerline{\normalsize{\faPhone\ 1xx-xxxx-xxxx \quad \faEnvelopeO\ \href{mailto:shlin0415@zju.edu.cn}{shlin0415@zju.edu.cn}}} 

\centerline{\normalsize{\faGithubSquare\ \href{https://github.com/shlin0415}{https://github.com/shlin0415} \quad \faExternalLink\ 2027届毕业生（杭州/远程）}} 
    
\section{\makebox[\widthof{\faGraduationCap}][c]{\color{CVBlue}\faGraduationCap}\ 教育背景}    
\textbf{浙江大学} \hfill 2023.9 -- 至今\\[0.2em] 
本科 \quad 浙江高考数学 135 / 英语 137
\begin{itemize}[nosep]
    \item 核心课程：数据结构与算法、计算机网络、操作系统、机器学习（均分：\textbf{填写你的高分课程或GPA}）
\end{itemize}

\section{\makebox[\widthof{\faUsers}][c]{\color{CVBlue}\faUsers}\ 实习与项目经历}

% --- 第一个项目：侧重AI Infra与Agent工程化 ---
\textbf{科研/医疗多智能体（Agent）系统开发} \hfill 华大基因 / M20 Genomics \titlebreak
项目描述：针对垂直领域（生物/医疗）构建具备自动化知识检索、实验设计能力的Agent，解决闭源API成本高及领域长文本推理延迟问题。
\begin{itemize}[nosep]
    \item \textbf{推理延迟优化}：通过引入\textbf{KV Cache预填充（Prefilling）与预加载}技术，在保持生成质量前提下，将语音/文本生成模型的首字延迟平均降低 \textbf{1.2s}。
    \item \textbf{高性能部署}：在 NVIDIA V100 等显卡上基于 \textbf{vLLM / sglang} 实现多模型量化部署，通过定制化的 \textbf{Sampling Kernel} 调查定位并修复了偶发性的生成低质（Bad Case）问题。
    \item \textbf{复杂系统设计}：构建动态模型路由，综合成本与隐私需求，实现云端 API 与本地低参数量模型（SDXL LoRA / Llama系列）的无感切换；负责特定领域 Knowledge/Rule 库的向量化封装与工具调用逻辑。
    \item \textbf{系统稳定性}：利用 \textbf{Snapshot} 分析排查 AI 推理框架中的 Memory Leak，通过持续跟踪 \textbf{lmcache} 等前沿技术优化长上下文（Long Context）检索效率。
\end{itemize}

\projectsep

% --- 第二个项目：侧重数据飞轮、算法与论文 ---
\textbf{PlantscRNAdb 4.0: 领域专用数据与自动标注系统} \hfill 核心开发者 / 一区论文发表 \titlebreak
项目描述：构建全球领先的单细胞转录组数据库（IF=24），旨在通过AI技术解决生物大数据标注成本高、一致性差的痛点。
\begin{itemize}[nosep]
    \item \textbf{数据飞轮构建}：设计高效的数据生产链路，通过对海量原始数据进行特征提取与噪声清洗，构建了高质量的训练/微调数据集。
    \item \textbf{自动标注优化}：优化现有标注算法模型，引入多模态融合特征，将领域内自动化标注准确率提升 \textbf{5\%+}，且推理速度优于现有流行方法，成功跑通“数据-模型-评估”的闭环。
    \item \textbf{成果}：相关成果发表于中科院一区期刊，负责数据库构建、算法优化及自动化流水线实现。
\end{itemize}

\section{\makebox[\widthof{\faCogs}][c]{\color{CVBlue}\faCogs}\ 技术栈与优势}
\begin{itemize}[nosep]
    \item \textbf{AI Infra \& 算法：} 精通 \textbf{vLLM, sglang, lmcache} 等推理加速框架；熟悉 \textbf{RLHF/PPO/DPO} 强化学习原理；深入理解 \textbf{RAG}、长短期记忆（Memory）及多智能体协作框架。
    \item \textbf{工程能力：} 熟练使用 \textbf{Python, Java (正在进修)}, Git, Linux 环境；具备 \textbf{CUDA Kernel} 级排查经验；能高效与 \textbf{Code Agent} 交互实现快速原型开发。
    \item \textbf{综合素质：} 具备极强的持续学习能力，深度关注 AI 技术前沿；对医疗科研、生物信息及 ACGN 领域有深刻理解。
\end{itemize}

\section{\makebox[\widthof{\faInfo}][c]{\color{CVBlue}\faInfo}\ 个人补充}
\begin{itemize}[parsep=0.5ex]
    \item \textbf{开源贡献：} \href{https://github.com/shlin0415}{github.com/shlin0415}（包含多类模型部署脚本与 Agent 实验 demo）
    \item \textbf{状态：} 暑期实习积极寻找中，杭州线下或线上（线上可立即到岗）。
\end{itemize}
\end{document}
